\section{Branching Processes}

Branching Processes are also known as Galton--Watson processes (1874) or Bienaym\'{e} processes (1845).

\begin{definition}[Branching Processes]
    \label{def:branching-process}%
    A \emph{branching process} is defined as follows. Let the population initially have \(X_0 = 1\), in generation \(0\). Let the number of offspring produced by each individual be \(X_1\), where
    \[
        \Prob\pqty{X_1 = k} = g_k, \quad k = 0, 1, 2, \ldots.
    \]

    Each individual produces an independent number of offspring, with the same distributions as \(X_1\). Let \(\pqty{Y_{k, n} : k \geq 1, n \geq 0}\) be independent and identically distributed random variables with the same distribution of \(X_1\), and we define \(Y_{k, n}\) be the number of offspring the \(k\)th individual of generating \(n\) produces. We define
    \[
        X_{n + 1} = Y_{1, n} + Y_{2, n} + \cdots + Y_{X_n, n}
    \]
    is the population at generation \(\pqty{n + 1}\).
\end{definition}

Our questions are: what is the probability that this population goes extinct, and what is the concrete distribution behind \(X_i\)?

\subsection{Distribution behind the Generations}

\begin{theorem}[Expected Number of Individuals in each Generation]
    \label{thm:branching-process-expt-generation}%
    The expectation of the number of individuals in each generation is given by
    \[
        \Expt\bqty{X_n} = \pqty{\Expt\bqty{X_1}}^{n}.
    \]
\end{theorem}

\begin{proof}
    We proceed by induction. Note this is trivial for \(n = 0\).

    Then, we have
    \[
        \Expt\bqty{X_{n + 1}} = \Expt\bqty{\Expt\bqty{\Conditional{X_{n + 1}}{X_n}}}.
    \]

    Conditioning upon the event \(X_n = m\), we have
    \begin{align*}
        \Expt\bqty{\Conditional{X_{n + 1}}{X_{n} = m}} &= \Expt\bqty{Y_{1, n} + \cdots + Y_{m, n}} \\
        &= m \cdot \Expt\bqty{X_1}.
    \end{align*}

    Therefore,
    \[
        \Expt\bqty{\Conditional{X_{n + 1}}{X_n}} = X_n \cdot \Expt\bqty{X_1}.
    \]

    Hence,
    \[
        \Expt\bqty{X_{n + 1}} = \Expt\bqty{X_n \cdot \Expt\bqty{X_1}} = \Expt\bqty{X_1} \cdot \Expt\bqty{X_n},
    \]
    finishing the inductive step, as desired.
\end{proof}

\begin{theorem}[Probability Generating Function of each Generation]
    \label{thm:branching-processes-pgf}%
    Let \(G\pqty{z} = \Expt\bqty{z^{X_1}}\) and \(G_n\pqty{z} = \Expt\bqty{z^{X_n}}\). Then
    \[
        G_{n + 1} \pqty{z} = G\pqty{G_{n} \pqty{z}} = G_n\pqty{G\pqty{z}} = G\pqty{G\pqty{\cdots G\pqty{z} \cdots}}.
    \]  
\end{theorem}

\begin{proof}
    We show from definition. We have
    \begin{align*}
        G_{n + 1} \pqty{z} &= \Expt\bqty{z^{X_{n + 1}}} \\
        &= \Expt\bqty{\Expt\bqty{\Conditional{z^{X_{n + 1}}}{X_n}}} \\
        &= \Expt\bqty{\Expt\bqty{\Conditional{z^{Y_{1, n} + Y_{2, n} + \cdots + Y_{X_n, n}}}{X_n}}}.
    \end{align*}

    Conditioning on a particular value of \(X_n = m\), we have
    \begin{align*}
        \Expt\bqty{\Conditional{z^{Y_{1, n} + Y_{2, n} + \cdots + Y_{X_n, n}}}{X_n = m}} &= \Expt\bqty{z^{Y_{1, n} + Y_{2, n} + \cdots + Y_{m, n}}} \\
        &= \pqty{G\pqty{z}}^{m}.
    \end{align*}

    Therefore,
    \[
        G_{n + 1} \pqty{z} = \Expt\bqty{\pqty{G\pqty{z}}^{X_n}} = G_n \pqty{G\pqty{z}}.
    \]
\end{proof}

\subsection{Extinction Probability}

We set
\[
    q = \Prob\pqty{X_n = 0 \text{ for some } n \geq 0},
\]
and
\[
    q_n = \Prob\pqty{X_n = 0}.
\]

\begin{definition}[Extinction Probability]
    \label{def:exctinction-probability}%
    The value \(q\) above is defined as the \emph{extinction probability}.
\end{definition}

Since we have \(\Bqty{X_n = 0} \subseteq \Bqty{X_{n + 1} = 0}\), we have \(q_n\) is increasing and converging to \(q\), by the continuity of probability measure (see \autoref{prop:continuity-of-probability-measure}).

\begin{proposition}[Extinction Probability is a Fixed Point]
    We have
    \[
        q_{n + 1} = G\pqty{q_n}, \qand q = G\pqty{q},
    \]
    where \(G\) is the probability generating function for \(X_1\).
\end{proposition}

\begin{proof}[of \(q = G\pqty{q}\) using \(q_{n + 1} = G\pqty{q_n}\)]
    Since \(G\) is continuous and \(q_n\) converges to \(q\), we have \(q = G\pqty{q}\).
\end{proof}

\begin{proof}[of \(q_{n + 1} = G\pqty{q_n}\), using Probability Generating Function]
    We have
    \begin{align*}
        q_{n + 1} &= \Prob\pqty{X_{n + 1} = 0} \\
        &= G_{n + 1}\pqty{0} \\
        &= G\pqty{G_{n} \pqty{0}} \\
        &= G\pqty{\Prob\pqty{X_n} = 0} \\
        &= G\pqty{q_n}.
    \end{align*}
\end{proof}

\begin{proof}[of \(q_{n + 1} = G\pqty{q_n}\), conditioning upon \(X_1\)]
    Conditioning upon \(X_1 = m\), let \(X_n^{\pqty{1}}, \ldots, X_n^{\pqty{m}}\) be independent and identically distributed branching processes with the same offspring as the original one, i.e., \(X_1\).

    Then we have
    \[
        X_{n + 1} = X_{n}^{\pqty{1}} + \cdots + X_n^{\pqty{m}}.
    \]

    Therefore,
    \begin{align*}
        q_{n + 1} &= \Prob\pqty{X_{n + 1} = 0} \\
        &= \sum_{m} \Prob\pqty{\Conditional{X_{n + 1} = 0}{X_1 = m}} \cdot \Prob\pqty{X_1 = m} \\
        &= \sum_{m} \Prob\pqty{\Conditional{X_n^{\pqty{1}} + \cdots + X_n^{\pqty{m}} = 0}{X_1 = m}} \cdot \Prob\pqty{X_1 = m} \\
        &= \sum_{m} \frac{\pqty{\Prob\pqty{X_n^{\pqty{1}} = 0}^{m}}}{q_n} \cdot \Prob\pqty{X_1 = m} \\
        &= G\pqty{q_n},
    \end{align*}
    as desired.
\end{proof}

The expectation and the extinction probability are linked in the following way.

\begin{itemize}
    \item If \(\Expt\bqty{X_1} < 1\), then we have \(\Expt\bqty{X_n} \to 0\), so we would expect the extinction probability to be \(q = 1\).

    \item If \(\Expt\bqty{X_1} > 1\), then we have \(\Expt\bqty{X_n} \to \infty\), so we would expect \(q < 1\).
    
    \item If \(\Expt\bqty{X_1} = 1\), then we have \(\Expt\bqty{X_n} = 1\) for all \(n\), but the value of \(q\) could be \(0\) (the trivial case) or \(1\) (everything else).
\end{itemize}

We want to find the extinction probability, so we solve the equation \(x = G\pqty{x}\). Note that \(x = 1\) is trivially a solution of this -- but we do not want that \(q = 1\), always.

\begin{figure}[ht!]
    \centering

    \begin{tikzpicture}[
        declare function={ f(\x) = 0.1 + 0.6 * \x + 0.3 * \x * \x; }
    ]
        \begin{axis}[
            axis lines = middle,
            xlabel = {\(x\)},
            ylabel = {\(y\)},
            xmin = -0.25, xmax = 1.25,
            ymin = -0.25, ymax = 1.25,
            xtick = {1},
            ytick = {1},
            domain = 0:1.05,
            samples = 100,
            clip = false
        ]
            \pgfmathsetmacro{\qa}{0}
            \pgfmathsetmacro{\qb}{f(\qa)}
            \pgfmathsetmacro{\qc}{f(\qb)}
            \pgfmathsetmacro{\qd}{f(\qc)}
            \pgfmathsetmacro{\qe}{f(\qd)}

            \pgfmathsetmacro{\q}{1/3}
            \pgfmathsetmacro{\fq}{f(\q)}

            \pgfmathsetmacro{\gf}{1.2}

            \addplot[only marks, mark=*, mark size = 1pt] coordinates {(0, 0) (1, 1) (\q, \fq) };

            \addplot[loosely dashed] coordinates {(1, 0) (1, 1) (0, 1)};

            \addplot[densely dotted] coordinates {(\qa, \qb) (\qb, \qb) (\qb, \qc) (\qc, \qc) (\qc, \qd) (\qd, \qd) (\qd, \qe)};

            \addplot[no marks] {x};
            \addplot[no marks] {f(x)};

            \addplot[dashed, domain = 1/6 : 1.1] {\gf * (x - 1) + 1};

        \end{axis}
    \end{tikzpicture}
    \begin{tikzpicture}[
        declare function={ f(\x) = 0.8 + 0 * \x + 0.2 * \x * \x; }
    ]
        \begin{axis}[
            axis lines = middle,
            xlabel = {\(x\)},
            ylabel = {\(y\)},
            xmin = -0.25, xmax = 1.25,
            ymin = -0.25, ymax = 1.25,
            xtick = {1},
            ytick = {1},
            domain = 0:1.05,
            samples = 100,
            clip = false
        ]
            \pgfmathsetmacro{\qa}{0}
            \pgfmathsetmacro{\qb}{f(\qa)}
            \pgfmathsetmacro{\qc}{f(\qb)}
            \pgfmathsetmacro{\qd}{f(\qc)}
            \pgfmathsetmacro{\qe}{f(\qd)}

            \pgfmathsetmacro{\gf}{0.4}

            \addplot[only marks, mark=*, mark size = 1pt] coordinates {(0, 0) (1, 1) };

            \addplot[loosely dashed] coordinates {(1, 0) (1, 1) (0, 1)};

            \addplot[densely dotted] coordinates {(\qa, \qb) (\qb, \qb) (\qb, \qc) (\qc, \qc) (\qc, \qd) (\qd, \qd) (\qd, \qe)};

            \addplot[no marks] {x};
            \addplot[no marks] {f(x)};

            \addplot[dashed, domain = -0.1 : 1.1] {\gf * (x - 1) + 1};

        \end{axis}
    \end{tikzpicture}

    \caption{Graph of \(G\pqty{x}\) and \(x\)}
    \label{fig:extinction-probability}
\end{figure}

We refer to \autoref{fig:extinction-probability} Note that \(q_0 = \Prob\pqty{X_0 = 0} = 0\), \(q_1 = G\pqty{q_0}\), and that \(q_n = G\pqty{q_{n - 1}}\). As this iteration continues, we see that \(q\) approaches the first non-negative intersection of the two graphs.

In the other case, if \(G\) does not intersect \(x\) at any point before \(1\), then the extinction probability will converge to \(1\).

The gradient of the tangent of \(G\) at \(1\) is given by \(G'\pqty{1^-} = \Expt\bqty{X_1}\). Note that in the first case, the slope is steeper than \(y = x\), so \(\Expt\bqty{X_1} > 1\); and in the second case, the slope is shallower than \(y = x\), so \(\Expt\bqty{X_1} < 1\).

These can be summarised into \autoref{thm:extinction-probability}.

\begin{theorem}[Extinction Probability]
    \label{thm:extinction-probability}%
    Suppose \(\Prob\pqty{X_1 = 1} < 1\). The extinction probability \(q\) is the minimal non-negative solution to the equation \(x = G\pqty{x}\).

    Moreover, \(q < 1\) if and only if \(\Expt\bqty{X_1} > 1\).
\end{theorem}

\begin{proof}[of the Value of \(q\)]
    We know that \(q = G\pqty{q}\). Let \(t\) be the smallest solution of \(x = G\pqty{x}\), and we will show that \(t = q\).

    We will prove by induction that \(q_n \leq t\), and hence by taking \(n \to \infty\) we have \(q \leq t\). But since \(t\) is the minimal solution, we must have \(q = t\).

    First, \(q_0 \leq t\) is trivially true, since \(q_0 = 0\). Suppose now \(q_n \leq t\), and we need to show \(q_{n + 1} \leq t\).

    We have
    \[
        q_{n + 1} = G\pqty{q_n} \leq G\pqty{t} = t,
    \]
    since \(G\) is an increasing function. This finishes the proof.
\end{proof}

\begin{proof}[of the Relation with Expectation]
    Now, we show that \(q < 1\) if and only if \(\Expt\bqty{X_1} > 1\).

    If \(g_0 + g_1 = \Prob\pqty{X_1 \leq 1} = 1\), we have \(\Expt\bqty{X_1} = g_1 < 1\), so \(G\pqty{z} = g_0 + g_1 z = 1 - \Expt\bqty{X_1} + \Expt\bqty{X_1} \cdot z\).

    Since \(G\pqty{z} = z\), we have \(\pqty{1 - \Expt\bqty{X_1}} \cdot z = 1 - \Expt\bqty{X_1}\). Since \(\Expt\bqty{X_1} = g_1 < 1\), it follows that the only solution is \(z = 1\), so the extinction probability \(q = 1\).

    Suppose from now on, assume that \(g_1 < 1\) and \(g_0 + g_1 < 1\). We want to show \(q < 1\) if and only if \(\Expt\bqty{X_1} > 1\).

    Define \(H\pqty{z} = G\pqty{z} - z = \sum_{r = 0}^{\infty} g_r z^r - z\), then we have \(H\pqty{1} = 0\). First we want to show that \(H\pqty{x}\) has at most one root in \(\ccintv{0}{1}\). We have
    \[
        H''\pqty{z} = \sum_{r = 0}^{\infty} r \pqty{r - 1} g_r z^{r - 2} > 0
    \]
    in \(\oointv{0}{1}\), since \(g_0 + g_1 < 1\). So \(H'\pqty{z}\) is strictly increasing in \(\oointv{0}{1}\), and hence \(H\) can have at most one more root other than \(1\). (If not, say there are roots \(0 < z_1 < z_2 < 1\), then \(H' = 0\) somewhere between \(z_1\) and \(z_2\) by Rolle's Theorem.)

    There are now two cases.
    
    \begin{itemize}
        \item If \(H\) has no other root than \(1\), then we have \(q = 1\). We note
        \[
            H'\pqty{1^-} = G'\pqty{1^-} - 1 = \Expt\bqty{X_1} - 1.
        \]

        We have \(H\pqty{0} = G\pqty{0} = g_0 \geq 0\), and \(H\pqty{1} = 0\), so \(H\pqty{z} \geq 0\) with \(z \in \ccintv{0}{1}\). Hence,
        \[
            H'\pqty{1^-} = \lim_{z \to 1^-}  \frac{H\pqty{z} - H\pqty{1}}{z - 1} \leq 0
        \]
        so \(\Expt\bqty{X_1} < 1\).

        \item If \(H\) has one more root, \(r < 1\), then \(r\) has to be the extinction probability.
        
        We have \(H\pqty{r} = 0\) and \(H\pqty{1} = 0\), with \(r < 1\), so by Rolle's Theorem, we have some \(z \in \oointv{r}{1}\) with \(H'\pqty{z} = 0\). We have \(H'\pqty{1^-} = \Expt\bqty{X_1} - 1\). \(H'\) is strictly increasing, and hence
        \[
            H'\pqty{z} < H'\pqty{1^-}
        \]
        which gives \(0 < H'\pqty{1^-}\), so \(\Expt\bqty{X_1} > 1\).
    \end{itemize}
\end{proof}